% Multimodal Large Language Models (MLLMs) have achieved remarkable success across diverse cognitive tasks, but they still struggle with the physical intelligence required to reason about the dynamic real world.
Multimodal Large Language Models (MLLMs) have achieved remarkable success across diverse expert-level tasks, but they still struggle with fundamental spatial-temporal abilities that humans naturally develop through continuous observation and interaction with the environment.
% Recent benchmarks have recognized this gap and begun to evaluate visual spatial intelligence, but they predominantly focus on static scene attributes, restrict evaluation to low-level perception, or rely on exact quantitative metrics that may conflate numerical prediction accuracy with genuine reasoning ability.
Recent benchmarks have recognized this gap and begun to evaluate visual spatial intelligence. However, they predominantly focus on static scene attributes, restrict evaluation to low-level perception, or rely on exact quantitative metrics that may conflate numerical prediction accuracy with genuine reasoning ability.
To address these limitations, we introduce the \textbf{Vi}sual \textbf{S}patial-\textbf{T}emporal \textbf{R}easoning \textbf{Bench}mark (\textbf{ViSTR-Bench}), a novel evaluation suite designed to systematically assess whether MLLMs can perform qualitative reasoning from continuous visual cues in dynamic scenes.
Guided by the principles of \textit{temporal emphasis}, \textit{reasoning orientation}, and \textit{qualitative evaluation}, ViSTR-Bench establishes a comprehensive four-dimensional framework covering Motion Perception, Spatial Relations, Outcome Prediction, and Physical Dynamics.
The benchmark comprises 12 distinct subtasks and 1,093 high-quality video question-answer pairs spanning diverse tabletop, indoor, and outdoor scenarios.
% Extensive evaluations across a broad spectrum of state-of-the-art proprietary, open-source, and specialized spatial MLLMs reveal that, despite their strong general video understanding capabilities, current models still face substantial bottlenecks in complex spatial-temporal reasoning and fall significantly short of human performance.
Extensive evaluations of a broad spectrum of state-of-the-art proprietary, open-source, and specialized spatial MLLMs reveal that, despite their strong general video understanding capabilities, current models still face substantial bottlenecks in complex spatial-temporal reasoning and remain far below human performance.